\honest{Anthropomorphic Optimism}{Eliezer Yudkowsky}{2008}

Anthropomorphism, I say, is expecting a process whose causes are unlike a human brain's to do what your brain predicts.

My example is one I have used before: the group selectionists, biologists before 1966 who believed predators would restrain their breeding so as not to exhaust their prey. Later Michael Wade selected insect populations for small size in the laboratory. Did the insects restrain their breeding and live ``in quiet peace with enough food for all, as the group selectionists had envisioned?'' No: they ate eggs and larvae, ``especially female larvae.'' \nb{Wynne-Edwards described competition and death, not quiet peace; see ``The Tragedy of Group Selectionism.''}

Why did the biologists not think of that? A tribe facing scarcity would agree to have fewer children; nobody would propose eating each other's daughters. The brain searches only among solutions it ranks high, so it never generates such ideas. Evolution ranks them differently.

In my line of work, I add, when you tell people an AI will ``not necessarily'' work like them, they give a reason why it must. I quote no one.

The group selectionists, I say, ``erased what had been written above the bottom line of their argument, without erasing the actual bottom line, and wrote in new rationalizations.'' \nb{The post quotes nothing the three biologists wrote. Their reasoning is inferred from their conclusion. Wynne-Edwards's own account, in 1980, is that the idea came to him while he was reviewing David Lack's book on the regulation of animal numbers.} Why would scientists do this? Humans ``seem to have evolved'' an instinct for arguing that their preferred policy serves everyone. A paragraph later, ``we have evolved an instinct to honestly believe'' it. I have added no evidence in between.

Evolution cannot be argued with, so the group selectionists were embarrassed. ``There's a fairly heavy subtext here about Unfriendly AI.'' That is all I say about AI.

The point generalizes, I conclude: optimism is ranking outcomes by your own preferences and taking the top one as your prediction. ``Look at the cognitive history and it's optimism in, optimism out.'' \nb{The general warning is sound. For its own example, the post does not look at the cognitive history.}
